\documentclass[preview]{standalone}

\usepackage{amsmath}
\usepackage{amssymb}
\usepackage{stellar}
\usepackage{definitions}
\usepackage{bettelini}

\begin{document}

\id{probability-discrete-independence}
\genpage

\section{Independence of discrete random variables}

\begin{snippetdefinition}{independent-discrete-random-variables-definition}{Independent discrete random variables}
    Let \((\Omega,\powerset(\Omega),\mathbb P)\) be a
    \countableprobabilityspace, and let
    \(X_i\colon\Omega\fromto E_i\), with \(i\in\{1,\ldots,n\}\), be
    \discreterandomvariable[discrete random variables], where
    \(n\in\naturalnumbers^\exceptzero\). The random variables
    \(X_1,\ldots,X_n\) are \emph{independent} if, for every
    \(A_i\subseteq E_i\),
    \[
        \mathbb P(X_1\in A_1,\ldots,X_n\in A_n)
        =
        \prod_{i=1}^n\mathbb P(X_i\in A_i).
    \]
\end{snippetdefinition}

\begin{snippetdefinition}{iid-discrete-random-variables-definition}{Independent and identically distributed discrete random variables}
    Let \((\Omega,\powerset(\Omega),\mathbb P)\) be a
    \countableprobabilityspace, and let
    \(X_1,\ldots,X_n\) be \discreterandomvariable[discrete random variables]
    with the same codomain \(E\). They are \emph{independent and identically
    distributed} if they are \independentdiscreterandomvariables and
    \[
        X_i\eqdist X_j
    \]
    for every \(i,j\in\{1,\ldots,n\}\).
\end{snippetdefinition}

\begin{snippetproposition}{independent-random-variables-subfamilies}{Subfamilies of independent random variables are independent}
    Let \((\Omega,\powerset(\Omega),\mathbb P)\) be a
    \countableprobabilityspace, and let
    \(X_i\colon\Omega\fromto E_i\), with \(i\in\{1,\ldots,n\}\), be
    \independentdiscreterandomvariables. If
    \(1\leq i_1<\cdots<i_k\leq n\), then
    \(X_{i_1},\ldots,X_{i_k}\) are independent.
\end{snippetproposition}

\begin{snippetproof}{independent-random-variables-subfamilies-proof}{independent-random-variables-subfamilies}{Subfamilies of independent random variables are independent}
    Let \(A_{i_j}\subseteq E_{i_j}\) for every \(j\in\{1,\ldots,k\}\).
    In the definition of independence of \(X_1,\ldots,X_n\), choose
    \(A_\ell=E_\ell\) for every
    \(\ell\notin\{i_1,\ldots,i_k\}\). Since
    \(\mathbb P(X_\ell\in E_\ell)=1\), the factorization becomes
    \[
        \mathbb P(X_{i_1}\in A_{i_1},\ldots,X_{i_k}\in A_{i_k})
        =
        \prod_{j=1}^k\mathbb P(X_{i_j}\in A_{i_j}).
    \]
\end{snippetproof}

\begin{snippetproposition}{independent-random-variables-events}{Events generated by independent random variables}
    Let \((\Omega,\powerset(\Omega),\mathbb P)\) be a
    \countableprobabilityspace, and let
    \(X_i\colon\Omega\fromto E_i\), with \(i\in\{1,\ldots,n\}\), be
    \independentdiscreterandomvariables. For every choice
    \(A_i\subseteq E_i\), the \event[events]
    \[
        \{X_i\in A_i\},\qquad i\in\{1,\ldots,n\},
    \]
    are \independentevents.
\end{snippetproposition}

\begin{snippetproof}{independent-random-variables-events-proof}{independent-random-variables-events}{Events generated by independent random variables}
    Let \(I\subseteq\{1,\ldots,n\}\). By the subfamily result,
    \[
        \mathbb P\left(\bigcap_{i\in I}\{X_i\in A_i\}\right)
        =
        \prod_{i\in I}\mathbb P(X_i\in A_i).
    \]
    This is exactly the defining condition for the family of events to be
    independent.
\end{snippetproof}

\begin{snippetproposition}{conditional-atom-characterization-independent-discrete-variables}{Conditional characterization at atoms}
    Let \((\Omega,\powerset(\Omega),\mathbb P)\) be a
    \countableprobabilityspace, and let
    \(X\colon\Omega\fromto E\) and \(Y\colon\Omega\fromto F\) be
    \independentdiscreterandomvariables. If \(x\in E\) satisfies
    \(\mathbb P(X=x)>0\), then, for every \(y\in F\),
    \[
        \mathbb P(Y=y\mid X=x)=\mathbb P(Y=y).
    \]
\end{snippetproposition}

\begin{snippetproof}{conditional-atom-characterization-independent-discrete-variables-proof}{conditional-atom-characterization-independent-discrete-variables}{Conditional characterization at atoms}
    By \conditionalprobability and independence,
    \[
        \mathbb P(Y=y\mid X=x)
        =
        \frac{\mathbb P(Y=y,X=x)}{\mathbb P(X=x)}
        =
        \frac{\mathbb P(Y=y)\mathbb P(X=x)}{\mathbb P(X=x)}
        =
        \mathbb P(Y=y).
    \]
\end{snippetproof}

\section{Density factorization}

\begin{snippettheorem}{independent-discrete-random-variables-density-factorization}{Density factorization criterion}
    Let \((\Omega,\powerset(\Omega),\mathbb P)\) be a
    \countableprobabilityspace, and let
    \(X_i\colon\Omega\fromto E_i\), with \(i\in\{1,\ldots,n\}\), be
    \discreterandomvariable[discrete random variables]. Let
    \(p_{X_1,\ldots,X_n}\) be their \jointdensity and let \(p_i\) be the
    \randomvariabledensity of \(X_i\). Then \(X_1,\ldots,X_n\) are
    independent if and only if, for every
    \((x_1,\ldots,x_n)\in\prod_{i=1}^n E_i\),
    \[
        p_{X_1,\ldots,X_n}(x_1,\ldots,x_n)
        =
        \prod_{i=1}^n p_i(x_i).
    \]
\end{snippettheorem}

\begin{snippetproof}{independent-discrete-random-variables-density-factorization-proof}{independent-discrete-random-variables-density-factorization}{Density factorization criterion}
    Suppose first that \(X_1,\ldots,X_n\) are independent. For every
    \((x_1,\ldots,x_n)\in\prod_{i=1}^n E_i\),
    \[
        p_{X_1,\ldots,X_n}(x_1,\ldots,x_n)
        =
        \mathbb P(X_1=x_1,\ldots,X_n=x_n)
        =
        \prod_{i=1}^n\mathbb P(X_i=x_i)
        =
        \prod_{i=1}^n p_i(x_i).
    \]

    Conversely, assume the displayed factorization. For every
    \(A_i\subseteq E_i\),
    \begin{align*}
        \mathbb P(X_1\in A_1,\ldots,X_n\in A_n)
        &=
        \sum_{x_1\in A_1}\cdots\sum_{x_n\in A_n}
        p_{X_1,\ldots,X_n}(x_1,\ldots,x_n)
        \\
        &=
        \sum_{x_1\in A_1}\cdots\sum_{x_n\in A_n}
        \prod_{i=1}^n p_i(x_i)
        \\
        &=
        \prod_{i=1}^n\sum_{x_i\in A_i}p_i(x_i)
        =
        \prod_{i=1}^n\mathbb P(X_i\in A_i).
    \end{align*}
    Therefore \(X_1,\ldots,X_n\) are independent.
\end{snippetproof}

\begin{snippetproposition}{two-variable-density-product-criterion}{Product form implies independence}
    Let \((\Omega,\powerset(\Omega),\mathbb P)\) be a
    \countableprobabilityspace, and let \(X\colon\Omega\fromto E\) and
    \(Y\colon\Omega\fromto F\) be \discreterandomvariable[discrete random
    variables]. Suppose that there exist \function[functions]
    \[
        \alpha\colon E\fromto[0,+\infty),
        \qquad
        \beta\colon F\fromto[0,+\infty)
    \]
    such that
    \[
        p_{X,Y}(x,y)=\alpha(x)\beta(y)
    \]
    for every \((x,y)\in E\cartesianprod F\). Then \(X\) and \(Y\) are
    independent.
\end{snippetproposition}

\begin{snippetproof}{two-variable-density-product-criterion-proof}{two-variable-density-product-criterion}{Product form implies independence}
    Set
    \[
        C=\sum_{y\in F}\beta(y),
        \qquad
        D=\sum_{x\in E}\alpha(x).
    \]
    Since \(p_{X,Y}\) is a \discretedensity,
    \[
        1=\sum_{x\in E}\sum_{y\in F}\alpha(x)\beta(y)=DC.
    \]
    In particular, \(C,D>0\). For every \(x\in E\) and \(y\in F\),
    \[
        p_X(x)=\sum_{y\in F}p_{X,Y}(x,y)=C\alpha(x),
        \qquad
        p_Y(y)=\sum_{x\in E}p_{X,Y}(x,y)=D\beta(y).
    \]
    Hence
    \[
        p_X(x)p_Y(y)=CD\alpha(x)\beta(y)=\alpha(x)\beta(y)=p_{X,Y}(x,y).
    \]
    The \densityfactorizationcriterion gives independence.
\end{snippetproof}

\section{Stability of independence}

\begin{snippetproposition}{joint-equality-in-law-preserves-independence}{Joint equality in law preserves independence}
    Let
    \((\Omega,\powerset(\Omega),\mathbb P)\) and
    \((\Omega',\powerset(\Omega'),\mathbb P')\) be
    \countableprobabilityspace[countable probability spaces]. Let
    \(X_i\colon\Omega\fromto E_i\) and
    \(Y_i\colon\Omega'\fromto E_i\), with \(i\in\{1,\ldots,n\}\), be
    \discreterandomvariable[discrete random variables]. If
    \[
        (X_1,\ldots,X_n)\eqdist(Y_1,\ldots,Y_n),
    \]
    then \(X_1,\ldots,X_n\) are independent if and only if
    \(Y_1,\ldots,Y_n\) are independent.
\end{snippetproposition}

\begin{snippetproof}{joint-equality-in-law-preserves-independence-proof}{joint-equality-in-law-preserves-independence}{Joint equality in law preserves independence}
    Equality in law of the two random vectors means that their joint densities
    agree. By marginalization, the one-dimensional marginal densities also
    agree. Therefore the joint density of \((X_1,\ldots,X_n)\) factorizes into
    its marginals if and only if the joint density of \((Y_1,\ldots,Y_n)\)
    factorizes into its marginals. The result follows from the
    \densityfactorizationcriterion.
\end{snippetproof}

\begin{snippetproposition}{functions-of-independent-random-variables}{Functions of independent random variables are independent}
    Let \((\Omega,\powerset(\Omega),\mathbb P)\) be a
    \countableprobabilityspace, and let
    \(X_i\colon\Omega\fromto E_i\), with \(i\in\{1,\ldots,n\}\), be
    \independentdiscreterandomvariables. For every \(i\), let
    \(f_i\colon E_i\fromto F_i\) be a \function, where
    \(F_i\subseteq\realnumbers^{d_i}\) is a finite or countably infinite
    \set for some \(d_i\in\naturalnumbers^\exceptzero\). Then
    \[
        f_1\circ X_1,\ldots,f_n\circ X_n
    \]
    are independent \discreterandomvariable[discrete random variables].
\end{snippetproposition}

\begin{snippetproof}{functions-of-independent-random-variables-proof}{functions-of-independent-random-variables}{Functions of independent random variables are independent}
    Let \(B_i\subseteq F_i\). Then
    \[
        \{f_i\circ X_i\in B_i\}
        =
        \{X_i\in f_i^{-1}(B_i)\}.
    \]
    Since \(X_1,\ldots,X_n\) are independent,
    \begin{align*}
        \mathbb P(f_1\circ X_1\in B_1,\ldots,f_n\circ X_n\in B_n)
        &=
        \mathbb P(X_1\in f_1^{-1}(B_1),\ldots,X_n\in f_n^{-1}(B_n))
        \\
        &=
        \prod_{i=1}^n\mathbb P(X_i\in f_i^{-1}(B_i))
        =
        \prod_{i=1}^n\mathbb P(f_i\circ X_i\in B_i).
    \end{align*}
\end{snippetproof}

\begin{snippetproposition}{independent-random-variables-disjoint-subvectors}{Disjoint subvectors of an independent family are independent}
    Let \((\Omega,\powerset(\Omega),\mathbb P)\) be a
    \countableprobabilityspace, and let
    \(X_1,\ldots,X_n\) be \independentdiscreterandomvariables. Let
    \(I,J\subseteq\{1,\ldots,n\}\) be disjoint nonempty \set[sets]. Then the
    \discreterandomvector[discrete random vectors]
    \[
        X_I=(X_i)_{i\in I},
        \qquad
        X_J=(X_j)_{j\in J}
    \]
    are independent.
\end{snippetproposition}

\begin{snippetproof}{independent-random-variables-disjoint-subvectors-proof}{independent-random-variables-disjoint-subvectors}{Disjoint subvectors of an independent family are independent}
    Since every subfamily of an independent family is independent, the family
    \((X_\ell)_{\ell\in I\union J}\) is independent. Its joint density
    factorizes into the product of its one-dimensional marginals. By
    marginalizing this factorization over the variables outside \(I\) and
    outside \(J\), the joint density of \((X_I,X_J)\) factorizes as the
    product of the marginal density of \(X_I\) and the marginal density of
    \(X_J\). The \densityfactorizationcriterion gives the independence of
    \(X_I\) and \(X_J\).
\end{snippetproof}

\begin{snippetproposition}{independent-random-variables-disjoint-sums}{Sums of disjoint independent subfamilies are independent}
    Let \((\Omega,\powerset(\Omega),\mathbb P)\) be a
    \countableprobabilityspace, and let
    \(X_1,\ldots,X_n\) be real-valued
    \independentdiscreterandomvariables. Let
    \(I,J\subseteq\{1,\ldots,n\}\) be disjoint nonempty finite
    \set[sets]. Then
    \[
        S_I=\sum_{i\in I}X_i,
        \qquad
        S_J=\sum_{j\in J}X_j
    \]
    are independent \discreterandomvariable[discrete random variables].
\end{snippetproposition}

\begin{snippetproof}{independent-random-variables-disjoint-sums-proof}{independent-random-variables-disjoint-sums}{Sums of disjoint independent subfamilies are independent}
    By the previous proposition, the subvectors \(X_I\) and \(X_J\) are
    independent. The sums \(S_I\) and \(S_J\) are functions of \(X_I\) and
    \(X_J\), respectively. Hence they are independent by the stability of
    independence under functions.
\end{snippetproof}

\end{document}
